\begin{scanpage}{0185}{167}
\[
\begin{aligned}E_\varphi=-H_\vartheta&=\frac{ik}{2\pi}H_{0x}\cos\varphi\frac{e^{ikR}}R\int n_ze^{i\Phi}\,dS,\\
E_\vartheta=H_\varphi&=\frac{ik}{2\pi}H_{0x}\frac{e^{ikR}}R\int(n_y\sin\vartheta+n_z\sin\varphi\cos\vartheta)e^{i\Phi}\,dS.
\end{aligned}\tag{26.11}
\]
Here $R$, $\vartheta$, $\phi$ are the spherical coordinates of the observation point, $\Phi=\psi-kr'\cos\Omega$, and integration is carried out over the illuminated elements of the body's surface. In the case of radar when the observation and irradiation directions coincide ($\vartheta=\pi-\gamma$, $\varphi=-\pi/2$), Equations (26.10) and (26.11) yield
\[
\left.\begin{aligned}E_x=-H_\vartheta&=-\frac{ik}{2\pi}E_{0x}\frac{e^{ikR}}R\int(n_y\sin\gamma+n_z\cos\gamma)e^{i\Phi}\,dS,\\
E_\vartheta=H_x&=0
\end{aligned}\right\}\tag{26.12}
\]
and
\[
\left.\begin{aligned}E_\vartheta=H_x&=\frac{ik}{2\pi}H_{0x}\frac{e^{ikR}}R\int(n_y\sin\gamma+n_z\cos\gamma)e^{i\Phi}\,dS,\\
E_x=H_\vartheta&=0.
\end{aligned}\right\}\tag{26.13}
\]
Furthermore, assuming the incident wave amplitudes are specified by Equation (26.02), let us write Expressions (26.12) and (26.13) in the following way:
\[
\left.\begin{aligned}E_x=-H_\vartheta&=\frac{iaE_{0x}}2\frac{e^{i\pi/4}}{\sqrt2}\frac{e^{ikR}}R\overline\Sigma^0,\\
E_\vartheta=H_x&=\frac{iaE_{0x}}2\frac{e^{-i\pi/4}}{\sqrt2}\frac{e^{ikR}}R\Sigma^0,
\end{aligned}\right\}\tag{26.14}
\]
where
\[
\Sigma^0=-\overline\Sigma^0=\frac k{\pi a}\int(n_y\sin\gamma+n_z\cos\gamma)e^{i\Phi}\,dS.\tag{26.15}
\]
Now let us represent the angular functions of fringing field (26.03) in the form
\end{scanpage}
